\subsection*{Paired question images change some openings while copying persists}
Training on real questions and retrieved passages allowed a direct test of what determined the first answer phrase. We trained two 628,104-parameter writers with the same initial weights and target examples. One received each original question. For the other we reassigned the questions across training examples while retaining the corresponding passages and human answers. With the original held-out questions, the two writers matched the opening four glyphs in 144 and 153 of 571 answers, respectively. Rotating the questions at test time reduced the matched writer's count to 126, but changed the reassigned writer's count by only one (Fig.~\ref{fig:grounded-glyph-writer}).

The correctly paired writer matched 32 complete references and the reassigned writer matched 37. All of these answers were contiguous in their inputs. The question-blind first-sentence baseline remained higher at 66. Of the matched writer's 144 correct openings, 112 continued or ended incorrectly. In the two illustrated questions the writer emitted every glyph of the supplied passage and stopped. The first question asked which script occupied the right-hand column; the second asked how many seals remained. In both cases the glyphs were clear, and the response repeated the premise. These observations locate a difficulty in choosing and completing the requested content even when the complete source shapes can be reproduced.

When supplied with two complete candidate answers and a common evidence image, the paired writer preferred the requested answer for both questions in both passage orders on 25 of 285 pairs. The reassigned writer did so on one pair. Comparing sets of pictured glyphs in the question and each answer gave 231 correct pairs with no training. Allowing a comparison via the two passage regions gave 251. The scores therefore carried information about question pairing, but lexical overlap explained much of the contrast in this constructed task.
